\subsection*{Which fields the computation exhausts}

Let $D<0$ be a fundamental discriminant, let $\ell$ be an odd prime
coprime to $D$, and put $E=\mathbb Q(\sqrt D)$ and $B=E(\zeta_\ell)$.
Write $r=\dim_{\mathbb F_\ell}\operatorname{Cl}(E)[\ell]$.
The number of fields in the reflection family attached to $D$ is
$(\ell^r-1)/(\ell-1)$. This counts isomorphism classes, and its
meaning includes a specified Galois action.

To see the action, take a nonzero class $[\mathfrak a]$ of order
$\ell$, write $\mathfrak a^\ell=(\alpha)$ and put
$n=N\mathfrak a>0$. Then $\alpha\bar\alpha=n^\ell$.
Changing the ideal representative changes $\alpha$ by an $\ell$th
power, up to a unit. Every unit of $E$ is an $\ell$th power under the
stated hypotheses. Thus the virtual units
\[
 \{v\in E^\times:(v)\ \text{is an }\ell\text{th ideal power}\}/E^{\times\ell}
 \simeq \operatorname{Cl}(E)[\ell]
\]
identify ideal-class lines with radical lines. The radical line remains
nonzero and distinct from the other lines in $B^\times/B^{\times\ell}$
because restriction followed by norm raises classes to $[B:E]=\ell-1$.

Set $M=B(\beta)$ with $\beta^\ell=\alpha$. A generator $s$ of
$\operatorname{Gal}(M/B)$ sends $\beta$ to $\zeta_\ell\beta$.
For a primitive root $g$ modulo $\ell$ choose $u$ fixing $E,\beta$
and sending $\zeta_\ell$ to $\zeta_\ell^g$. Complex conjugation $c$
sends $\beta$ to $n/\beta$ and $\zeta_\ell$ to $\zeta_\ell^{-1}$.
These automorphisms generate the full Galois group. Direct substitution
shows that $c$ is central and $usu^{-1}=s^g$.
The $\ell$ real conjugates of $\beta+n/\beta$ are
$\zeta_\ell^j\beta+n/(\zeta_\ell^j\beta)$.
They are distinct. Equality for different $j$ would give
$\beta^2/n$ an $\ell$th root of unity and hence $\alpha^2=n^\ell$,
forcing $[\alpha]=0$ in $E^\times/E^{\times\ell}$.
Their splitting field is $M^{\langle c\rangle}$,
since the affine action of $s$ and $u$ on their indices is faithful.
Its cyclic degree-$(\ell-1)$ subfield $B^{\langle c\rangle}$ is the
mirror field in the sense of Cohen, Rubinstein-Salzedo and Thorne.
The displayed conjugation rule is exactly their star condition.

Two different ideal-class lines give different normal closures.
On adjoining $E$ each normal closure recovers its Kummer extension
$M/B$, and Kummer theory distinguishes the lines. Different generators
of one line give the same normal closure. Its cyclic complements of
order $\ell-1$ are conjugate in the affine group, so they give a single
isomorphism class of degree-$\ell$ fields.

For each of the discriminants used here, the local discriminant
calculation places these fields among the totally real starred fields
of discriminant $|D|^{(\ell-1)/2}\ell^{\ell-2}$ or
$|D|^{(\ell-1)/2}\ell^\ell$.
Theorem~1.2 and Lemma~2.8 of Cohen, Rubinstein-Salzedo and Thorne
give precisely $(\ell^r-1)/(\ell-1)$ starred fields in the union
of these two sets. The construction therefore exhausts this family.
The discriminant alternatives and the pooled count are established
in that paper. The present table assigns a specific alternative and
defining polynomial to each ideal-class line.
\href{https://arxiv.org/pdf/1405.1075v2}{Cohen, Rubinstein-Salzedo
and Thorne, Theorem~1.2, Lemma~2.8 and Proposition~3.10}.

The analytic input enters in exhausting the base fields.
Elsenhans, Kl\"uners and Nicolae, Theorem~1 and Table~1, find 27
imaginary quadratic fields of class-group exponent five through
$3.1\cdot 10^{20}$. Their Lemma~1 and Theorem~3 bound every possible
absolute discriminant of exponent five by $2.3\cdot 10^{20}$ under the
quadratic least-split-prime estimate (Table~2). The requisite critical-line
statement for $\zeta_E(s)=\zeta(s)L(s,\chi_D)$ follows from the public
Dirichlet GRH theorem. Their published data give 25 groups $C_5$ and
two groups $C_5^2$, at $D=-12451,-37363$. Thus the count is
$25+2(5^2-1)/(5-1)=37$.
\href{https://arxiv.org/pdf/1803.02056}{Elsenhans, Kl\"uners and
Nicolae, Theorems~1--3 and Tables~1--2}.
\href{https://math.uni-paderborn.de/en/ag/ca/research/exponent/exponent0}
{Authors' tables of discriminants and groups}.
